\section{Discussion}\label{sec:discussion}

\paragraph{What a practitioner can take away.} Table~\ref{tab:recipe} lists each component of the method with when to use it, its measured effect, where it does not help, and its cost; Table~\ref{tab:defaults} gives the shipped settings and Appendix~\ref{app:protocol} the prompts. Three quantities decide most choices at run time: the degeneracy flags decide which count to read; agreement between two calls of one model measures stability, not correctness, so high agreement means another such call adds little; and the family default, or $p_1$ estimated on a few labelled images, decides whether to delete unconfirmed records.

\begin{table*}[tp]
\iflatexml\else\centering\fi\footnotesize\setlength{\tabcolsep}{3.5pt}\renewcommand{\arraystretch}{1.15}
\caption{FSR at a glance: what each component is for, its measured effect, where it gives no significant gain or harms, and its cost. Effects are relative changes of the stated loss against one call of the same model unless stated; ``n.s.'': not significant. The last row lists remedies with small or no significant gains.}\label{tab:recipe}
\begin{tabular}{@{}>{\raggedright\arraybackslash}p{2.5cm}>{\raggedright\arraybackslash}p{2.6cm}>{\raggedright\arraybackslash}p{3.8cm}>{\raggedright\arraybackslash}p{4.0cm}>{\raggedright\arraybackslash}p{1.9cm}@{}}\toprule
Component & Use when & Measured effect & No significant gain, or harm, when & Cost \\\midrule
List first, read the list count; stated count only if the list is flagged & counting, above all several classes in dense scenes & per-type count error $-85.2\%$ (GPT-6.1), $-98.8\%$ (Sonnet 5.5); same response, readout only: $-71\%$ to $-84\%$ (\S\ref{sec:res:readout}) & one-class photographs, and aerial or document images whose lists miss objects: small, n.s.\ for GPT-6.1 and Sonnet 5.5 (\S\ref{sec:res:readout}) & one call \\
Allowed names in the prompt, aliases mapped in code & the model uses names outside the vocabulary & Sonnet 5.5's ``diamond'': error 5.42 to 0.42 on the same lists (Table~\ref{tab:readout}) & names already in the vocabulary & no extra calls \\
Degeneracy flags; Gemini 3.8 at its default setting for counts & long lists that may break & flagged lists 90\% to 12\%, count error $-53.6\%$ at the same cap (Table~\ref{tab:setting}) & Gemini 3.8 at its lowest setting, even with the fallback: $+15.6\%$, n.s. & more reasoning tokens \\
Two-stage zoom & clicks, small targets & miss rate $-65\%$ (Gemini 3.8), $-41\%$ (Sonnet 5.5) (\S\ref{sec:res:view}) & GPT-6.1: n.s. & about 2 calls \\
Views from a planner, merged; each view confirms only records it contains & small objects; views from a planner or crops & four crops (6\% overlap) confirming a whole-image report: set error up to $-45.5\%$, better in 4 of 6 cells (\S\ref{sec:res:view}) & where unconfirmed records are mostly true ($p_1>1/2$): worse in 2 of 6 cells; tiles without a one-object overlap lose large objects & one call per view \\
Three-call unanimity, coordinates averaged & set error on over-reporting families (detection, layout, screens) & significantly better in 9 of 15 cells, significantly worse in none (\S\ref{sec:res:transfer}) & F1 (2 better, 4 worse, 12 n.s.); under-reporting dense points (worse in 2 of 3) & 3 calls \\
Per-model single-call configuration for F1 & F1 on detection-type sources & macro F1 42.1 to 49.1; significantly better in 5 of 12 cells, worse in 1 (\S\ref{sec:res:transfer}) & set error: 2 better, 1 worse, 9 n.s. & $1.44\times$ the cost of the same model's list call at the lowest setting \\
Cross-model alarm & deciding which items to check & catches 61\% of GPT-6.1's misses against 17\% for a same-model call; 68\% to 78\% after zoom (\S\ref{sec:res:persist}) & flags more items; disagreement does not say which answer is right, and agreement does not certify & one call of a cheaper model \\
Three-model GUI-grounding tier (Algorithm~\ref{alg:clicks}) & higher click accuracy is worth up to 15 planned calls per item & miss rate 13.67\% to 9.33\%; 3.00 points better than a plain vote over the same points (\S\ref{sec:res:choose}) & against crop-first: n.s.; a fourth source: n.s.; the comparison step alone: n.s. & \$0.0273 per item, $3.2\times$ one zoom \\
Remedies with small or no significant gains & \multicolumn{4}{>{\raggedright\arraybackslash}p{12.7cm}}{a three-click vote or zoom gated on disagreement (miss-rate reductions of at most 13\%); grids, ticks and cell labels (localization changes $\le0.23$ px); ``do not miss any'' wording; votes over repeated calls of one model (error reductions of at most 23\% for counts and 4.5\% for clicks); same-view self-review on screens; an equal-weight vote of all models when one is strongest (\S\ref{sec:res:view}, \S\ref{sec:res:common}, \S\ref{sec:res:transfer})} \\\bottomrule
\end{tabular}
\end{table*}

\paragraph{Why the gains come from changing information.} The theory and the measurements point the same way. Object-level persistence means that repeated calls of one model share most of their errors. In our comparisons, votes over the same outputs gained at most 23\% for counts and 4.5\% for clicks, and even a perfect selector over three calls can only choose among the records those calls contain. The largest gains we measured came from operations that change what is available: reading a different field of the same response, looking at a different view, or asking a model with a different model-specific component. Where those operations fail, the failures have the predicted form: deletion hurts when $p_1>1/2$, tiling can lose objects that no tile contains whole, and a partner hurts when its additional loss of true matches outweighs its additional rejection of false records, with the effect of averaged coordinates counted separately.

\paragraph{Where click accuracy can still come from.} Of the tier's 56 remaining misses on the 600 click items, 43 have no correct candidate and 13 are lost in representation or choice (Proposition~\ref{prop:three}). A better rule over the same candidates can act only on the 13. The 43 need a source with a high rescue rate $\phi$ on exactly those items, and its gain must survive representation and choice: of the crop-first source's 7 rescues, 3 did (Section~\ref{sec:res:choose}). On items where the models agree, any further action must meet the break-rate condition of Proposition~\ref{prop:switch}. These two measurable quantities, the rescue rate on the items no candidate covers and the break rate on agreeing items, decide whether a further source or action raises accuracy, and both are measured from stored calls on labelled items.

\paragraph{What future models should change.} The measurements point to properties that model builders can target directly. A model's stated summary should be computed from its own enumeration: GPT-6.1 and Sonnet 5.5 already list almost exactly the right number of objects, and their stated counts lose that information. Long lists should not degrade with position; Gemini 3.8 at its lowest setting drops class fields more often when consecutive records share a class; an order that interleaves classes lowers its count error and grouping by class raises it (Section~\ref{sec:res:degen}). Class vocabularies given in the prompt should be respected, since Sonnet 5.5 still wrote ``diamond'' 258 times when the format named the allowed classes. Finally, errors that persist across calls limit what a fixed budget of resampling can repair; diversity across calls, or calibrated uncertainty at the object level, could make the cheaper same-model alarm as useful as a second model.

\paragraph{Scope of the results.} The theory's guarantees hold under their stated conditions. The identities of Section~\ref{sec:theory:identities} hold for any outputs and any fixed choice of maximum matching. Proposition~\ref{prop:partner} assumes deletion-only confirmation, no match takeover, and no ignore region or none of the base report's false records inside one; outside these conditions the exact identity of Proposition~\ref{prop:identity} replaces it. The disagreement floor is a bound for point tasks and for boxes at IoU thresholds above 0.5; the click alarm is empirical and is evaluated by risk and acceptance coverage. Proposition~\ref{prop:nocert} concerns model classes that admit scenes with identical or nearly identical output laws; when an object is reported with probability at most $\vartheta$ per call, $K$ calls separate the two scenes by at most $K\vartheta$ in total variation. The readout repair applies where a readout gap exists: on our dense synthetic scenes the gap is most of the error, while on aerial and document images it is a small part of it for GPT-6.1 and Sonnet 5.5. Section~\ref{sec:experiments} states how the reasoning setting and the output cap were separated. All results are for the model versions listed in Appendix~\ref{app:registry}, at a total cost of about \$758; we release every response and the analysis code, so the reported identities and tables can be reproduced without new model calls, while evaluating a new model version requires its own responses, analysed the same way.

\paragraph{Broader impact.} More accurate and auditable visual reports can make inventory, document and screen-automation systems more reliable, and an explicit status and alarm make failures visible instead of silent. The same techniques could make automated agents act with more confidence; our pipeline always returns a status and never claims certification without labels.
